\paragraph{Map an exact insurer schedule into a proposed book.}
A separate retained insurer calculation supplies the missing posting example.
It starts with cash of 200 and an insurance liability of 100, including a loss
component of 20. Its six journals already specify exact amounts; the new
adapter must show where those amounts go in the company book and what state
must accompany them. This is a different synthetic case from the resolved
assessment above. The original calculator continues to return a report without
posting, while \code{InsuranceLossBook/1} proposes a separate sequence of
declared company methods.

The request supplies the opening book, an injective mapping for six account
roles, three typed memory locations and a statement mapping. The account roles
are cash, insurance liability, finance expense, service revenue, loss-reversal
income and claims expense. The memory records the loss component, the next
step and the exact native request identity. Before constructing a method, the
adapter recalculates the selected native request and checks the complete result,
opening balances, account classes and memory. It accepts only the two retained
requests; a plausible new schedule is outside this checked version.

The illustrative units require an explicit bridge as well. The native case
uses \code{CU}, whereas the host book requires a three-letter currency field.
The request binds both complete contexts and declares \code{XCU} to denote
the same illustrative unit at scale one. No exchange-rate calculation occurs.
The entity remains Synthetic Insurer, and the native group and historical
rule stay attached to the original request. This syntactic alias does not
admit a currency, accounting basis or real entity.

The following table gives all six signed journals for the nonzero-loss case.
Amounts use the same illustrative units; debits are positive and credits
negative. Each date has finance, liability release and claim-payment entries,
in that order. A supplied premise classifies the claim payments as operating
cash flows.
\begin{center}\small
\begin{tabular}{lll}
\toprule
Date & Entry & Signed account changes \\
\midrule
1 Jan 2026 & Finance & Expense $+10$; liability $-10$ \\
           & Release & Liability $+55$; revenue $-44$; reversal $-11$ \\
           & Claims  & Claims expense $+50$; cash $-50$ \\
1 Jan 2027 & Finance & Expense $-5$; liability $+5$ \\
           & Release & Liability $+50$; revenue $-40$; reversal $-10$ \\
           & Claims  & Claims expense $+45$; cash $-45$ \\
\bottomrule
\end{tabular}
\end{center}
Every row balances without an opening-capital entry or a balancing amount.
After the first three entries cash is 150, liability is 55 and profit is
$44+11-10-50=-5$. The next three leave cash at 105 and liability at zero;
their profit is $40+10-(-5)-45=10$. The negative finance expense in the second
date is retained explicitly, rather than replaced by a positive expense.

Let $K$ denote the loss component measured in the same units as the liability.
It is tracked within that liability, not added to it as a separate obligation.
The retained calculation attributes its changes separately to finance, claim
release and risk-adjustment release. Under its supplied allocation assumptions,
the two date-end values are
\[
  K_1=20+2-10-1=11,\qquad
  K_2=11-1-9-1=0.
\]
The release method combines the two release amounts in one memory update but
keeps their original explanations separately. In the second retained case,
$K$ starts at zero and stays there. The liability releases are then revenue
of 55 and 50, with zero loss reversals. Cash, liability and total profit match
the first case, but the gross revenue and reversal accounts differ. Retaining
the zero postings and separate origins makes that difference inspectable.

Every generated event checks its original date, scoped opening balances,
loss component, progress and case identity before proposing its changes.
Postconditions check the resulting balances and memory; skipped or repeated
steps refuse. Unrelated equipment, equity and quoted or Unicode text in memory
survive unchanged. Thus a successful calculation does not require replacing
the whole company with the insurer example.

The host interval starts on 1 January 2025 and ends exclusively on
2 January 2027 so that both native event dates fit without alteration. Reports
after the third and sixth entries are cumulative within that interval: profit
is $-5$ and then $5$, and operating payments are 50 and then 95. These reports
do not close annual accounts or move 1 January events into the preceding year.
Annual closing and reopening need their own explicit treatment.

Complete replay covers both native cases on two account charts, with and
without an unrelated equipment balance. Unknown or false native findings,
changed contexts, forged results, invalid mappings and inconsistent state
prevent a valid proposal. The retained evidence is in
\path{data/accounting-master/insurance-loss-book-20261025/}.
It establishes this exact mapping and proposed execution. Source applicability,
broader insurance measurements, durable posting and a complete statement-note
package remain separate requirements; the supplied historical findings and
allocation assumptions are preserved as such.
